	\section{Módulos de la Plataforma}
	\label{cap:modulos}
	Este capítulo recorre los módulos entregados desde la perspectiva de quien los usa. Mientras el
	Capítulo~\ref{cap:implementacion} describe cómo se construyó el sistema por capas técnicas, aquí
	se describe qué hace cada aplicación, pantalla por pantalla, y cómo esas pantallas se
	corresponden con los casos de uso definidos en la Sección~\ref{sec:casos_uso}. El orden sigue el
	recorrido natural de la información dentro de la institución: entra por el portal del quejoso,
	se procesa en el panel de revisión y se sostiene desde la consola de administración.
	
	\subsection{Portal del quejoso}
	\label{sec:modulo_quejoso}
	El portal del quejoso es la única aplicación de acceso público y, por lo tanto, la que define la
	experiencia que la comunidad politécnica tiene de la Defensoría en línea. Su diseño parte de una
	premisa que condicionó todo lo demás: la persona que necesita presentar una queja no debería
	tener que crear una cuenta antes de poder hacerlo.
	
	\subsubsection{Área pública}
	El área pública es accesible sin credenciales y agrupa cuatro funciones.
	
	La \textbf{página de inicio} presenta la identidad institucional y tres accesos directos:
	presentar una queja, consultar el estado de una queja existente e iniciar sesión. Incorpora
	además una sección de servicios en línea que replica los accesos de referencia del sitio
	institucional de la Defensoría, de modo que el portal se perciba como continuación del servicio
	existente y no como un sistema desconectado.
	
	El \textbf{formulario de registro de queja} implementa el caso de uso CU-Q01 y es la pantalla más
	compleja del portal. Está organizada en secciones numeradas: datos del quejoso, datos de la
	queja, evidencias y, cuando corresponde, datos del tutor. Tres decisiones de diseño merecen
	mención. La unidad académica donde ocurrieron los hechos se selecciona del catálogo institucional
	en lugar de capturarse como texto libre, lo que hace el dato comparable entre expedientes. Las
	evidencias admiten varios archivos con lista previa de lo seleccionado y opción de retirar
	alguno antes de enviar. Y el aviso normativo sobre el plazo para presentar una queja se presenta
	en un elemento flotante consultable en cualquier momento, en lugar de un bloque fijo que
	compitiera por el espacio del formulario.
	
	El \textbf{tratamiento de menores de edad} se resuelve dentro del mismo formulario. Cuando la
	fecha de nacimiento indica que el quejoso es menor, el sistema exige los datos de un tutor o
	adulto responsable antes de permitir el envío, conforme a la regla RN-Q01. Una vez capturados, el
	formulario muestra un aviso persistente que confirma que ya se registraron y ofrece editarlos,
	porque en una versión previa no existía ninguna señal visual de que ese paso se hubiera
	completado y los usuarios volvían a abrir el diálogo por duda.
	
	La \textbf{consulta por folio} implementa el caso de uso CU-Q02 y permite conocer el estado de un
	expediente con el folio y el correo con que se presentó, sin necesidad de cuenta. Es la
	funcionalidad que resuelve de forma más directa el problema de trazabilidad que motivó el
	proyecto.
	
	Completa el área pública un \textbf{menú de preguntas frecuentes} sobre la Defensoría y el
	proceso de queja, cuyo contenido administra el personal de sistemas sin intervención del equipo
	de desarrollo.
	
	\subsubsection{Panel autenticado}
	Tras activar una cuenta con el folio de una queja previa, el quejoso accede a un panel con seis
	secciones.
	
	\begin{itemize}
		\item \textbf{Resumen.} Conteo de las quejas propias agrupadas por estatus, como vista de
		entrada al panel.
	
		\item \textbf{Mis Quejas.} Listado con filtros por estatus y por texto, y detalle de cada
		expediente con su línea de tiempo. Desde el detalle se editan los campos habilitados, pero
		únicamente mientras la queja no haya iniciado su validación: al pasar a validación los campos
		se bloquean, y el servidor rechaza la modificación aunque se intente por otra vía.
	
		\item \textbf{Nueva Queja.} Registro de una queja adicional sin volver a capturar la
		identidad, que se toma de la sesión.
	
		\item \textbf{Acuerdos de Conciliación.} Consulta de los acuerdos propuestos por la
		Defensoría, con la posibilidad de aceptarlos o rechazarlos y añadir un comentario.
	
		\item \textbf{Notificaciones.} Historial persistente de avisos ---accesos, quejas creadas,
		cambios de estatus, acuerdos--- con contador de no leídas que se actualiza al marcarlas.
	
		\item \textbf{Perfil.} Edición de los datos de contacto. El nombre, el correo institucional y
		la boleta o número de empleado se muestran en modo de solo lectura, porque provienen de la
		activación de la cuenta y modificarlos rompería la correspondencia con las quejas ya
		presentadas.
	\end{itemize}
	
	\subsection{Panel de revisión}
	\label{sec:modulo_revision}
	El panel de revisión es la herramienta de trabajo diario del personal de recepción y el punto
	donde una queja deja de ser un formulario recibido para convertirse en un expediente en trámite.
	Su diseño se organiza alrededor de una secuencia fija: revisar, decidir, registrar.
	
	\subsubsection{Bandeja de entrada}
	La bandeja es la pantalla de inicio y presenta contadores de quejas pendientes, en proceso y
	turnadas en el día, junto con la lista de expedientes por trabajar. La elección de encabezar con
	contadores responde a una necesidad expresada en las sesiones de levantamiento: el personal
	necesita saber de un vistazo cuánta carga tiene antes de empezar a abrir expedientes.
	
	\subsubsection{Detalle y validación}
	Al abrir un expediente, el personal accede al resumen completo de la queja y a sus documentos
	adjuntos, que puede previsualizar o descargar. La pantalla incorpora además una \textbf{búsqueda
	de antecedentes}: otras quejas previas de la misma persona. Esta función atiende un problema real
	del proceso, la duplicidad de expedientes cuando alguien presenta la misma queja por distintos
	canales, y su ubicación en la pantalla de decisión es deliberada: el dato debe estar disponible
	en el momento de decidir, no en un módulo separado que haya que consultar por iniciativa propia.
	
	Desde esta pantalla el personal toma una de dos decisiones excluyentes.
	
	El \textbf{rechazo} exige seleccionar al menos un motivo y admite observaciones libres. Al
	confirmarlo, el expediente pasa a estado rechazado y el sistema envía de forma automática el
	correo correspondiente al quejoso. La regla RN-R03 hace que el rechazo nunca sea una acción
	silenciosa: no es posible descartar un expediente sin que su autor se entere ni sin dejar
	registrado el motivo.
	
	El \textbf{turnado} exige elegir el área o dependencia ---tomada del catálogo institucional--- y
	el defensor o subdefensor responsable, tomado del catálogo de personal activo con esos roles.
	Admite comentarios y genera el folio oficial de turnado, distinto del folio de seguimiento que
	conoce el quejoso.
	
	En ambos casos el sistema registra quién procesó el expediente y cuándo, sin excepción, conforme
	al requerimiento RNF-R01.
	
	\subsubsection{Registro manual de quejas en papel}
	Una parte del volumen que recibe la Defensoría llega en papel. El módulo de registro manual
	permite incorporar esos expedientes al mismo flujo digital, capturando la identidad del quejoso,
	la dependencia, el número de oficio, la fecha de recepción física, el tipo de documento, la
	descripción, la ubicación física del expediente y, de forma opcional, el archivo digitalizado.
	
	Este módulo tiene una particularidad de diseño: admite que no se conozca la boleta o el número de
	empleado del quejoso, situación que el registro propio no permite. La razón es que un documento
	en papel puede provenir de una persona externa o llegar sin esos datos, y rechazarlo por ello
	significaría dejar fuera del sistema quejas que la institución sí recibió. El expediente queda
	marcado con origen manual, de modo que su procedencia siempre es identificable.
	
	\subsubsection{Historial y exportación}
	El historial reúne los expedientes ya procesados ---turnados o rechazados--- con filtros por
	texto, estatus y fecha. Su función principal no es la consulta puntual, que ya resuelve la
	búsqueda de antecedentes, sino la generación de reportes: el resultado filtrado se exporta a un
	archivo de hoja de cálculo que se abre directamente en las herramientas ofimáticas que el
	personal ya usa, sin conversión previa.
	
	\subsubsection{Emisión de acuerdos de conciliación}
	El personal puede emitir un acuerdo de conciliación dirigido al quejoso de un expediente,
	capturando asunto y términos. El acuerdo queda en estado pendiente y aparece de inmediato en el
	panel del quejoso, que puede aceptarlo o rechazarlo. Este módulo es el punto donde el flujo de la
	plataforma se conecta con las etapas del procedimiento que aún no tienen interfaz propia.
	
	\subsection{Consola de administración}
	\label{sec:modulo_admin}
	La consola de administración es la que hace que la plataforma sea operable a largo plazo sin
	intervención del equipo de desarrollo. Agrupa cuatro frentes independientes entre sí.
	
	\subsubsection{Usuarios y roles}
	Ciclo de vida completo de las cuentas del personal: alta con contraseña temporal generada por el
	sistema, edición, restablecimiento de contraseña, baja y reactivación, para cualquiera de los
	cinco roles.
	
	Dos decisiones de diseño de este módulo tienen consecuencias de seguridad. La primera es que las
	cuentas nuevas se crean siempre como temporales y con obligación de cambiar la contraseña en el
	primer acceso, de modo que la contraseña generada por el sistema nunca sea la contraseña
	permanente de nadie. La segunda es que dar de baja una cuenta la desactiva en lugar de
	eliminarla: si se eliminara, las acciones que esa cuenta registró en la bitácora quedarían sin
	autor identificable, lo que anularía el propósito de la auditoría.
	
	\subsubsection{Catálogo de unidades académicas}
	Administración del catálogo institucional que consumen el formulario público de quejas y el
	turnado del panel de revisión. Permite el alta y la edición individual, y la importación masiva
	desde un archivo de hoja de cálculo con el formato institucional.
	
	La importación masiva se diseñó para ser tolerante a errores: en lugar de rechazar el archivo
	completo cuando una fila está mal formada, procesa lo que puede y devuelve un resumen de filas
	procesadas, creadas, actualizadas y con error. El criterio detrás de esa decisión es que un
	archivo de doscientos registros con un error tipográfico en uno de ellos no debería obligar a
	repetir toda la carga.
	
	\subsubsection{Plantillas de documentos oficiales}
	Administración del contenido reutilizable de los oficios que genera el procedimiento. Cada
	plantilla admite marcadores que el sistema sustituye por los datos del expediente al generar el
	documento, y la consola permite consultar qué marcadores están disponibles, previsualizar el
	resultado con datos de ejemplo y publicar los cambios.
	
	Cada plantilla se identifica por una clave estable de negocio y no por su identificador
	autogenerado, de modo que el resto del sistema pueda referenciarla sin depender del orden en que
	se hayan creado. Cada actualización registra quién la hizo y cuándo.
	
	\subsubsection{Seguridad y respaldos}
	Este módulo agrupa la operación más sensible de la plataforma. Permite listar los respaldos
	existentes, ejecutar un respaldo bajo demanda, descargar uno específico y restaurar la base de
	datos a partir de él.
	
	La restauración es la única operación destructiva de todo el sistema y recibe un tratamiento
	diferenciado: exige una confirmación explícita en el cuerpo de la petición, de modo que no pueda
	dispararse por accidente con una sola acción de la interfaz. Además, la capacidad de generar,
	descargar y restaurar respaldos está reservada al rol de administración de sistemas, que es el
	único que dispone de esos puntos de acceso.
	
	El módulo incluye también la \textbf{bitácora de acciones críticas}, que registra las altas,
	bajas, ediciones y reactivaciones de cuentas, los restablecimientos de contraseña, las
	actualizaciones de plantillas, los respaldos manuales, las restauraciones y los inicios de
	sesión, con el usuario que los ejecutó, la dirección de origen y la fecha. Esta bitácora es el
	mecanismo que hace verificable el requerimiento RNF-A02 y, junto con la regla de que las cuentas
	se desactivan en lugar de eliminarse, constituye el soporte de auditoría de la plataforma.
	
	\subsubsection{Administración de preguntas frecuentes}
	Alta, edición y eliminación del contenido del menú de preguntas del portal público. Los cambios
	se reflejan de inmediato en el portal, sin necesidad de recompilar ni redesplegar la aplicación
	del quejoso, lo que traslada el mantenimiento de ese contenido del equipo de desarrollo al
	personal de la institución.
	
	\subsection{Relación entre módulos y casos de uso}
	\label{sec:trazabilidad_modulos}
	La Tabla~\ref{tab:trazabilidad} presenta la correspondencia entre los módulos entregados y los
	casos de uso especificados en la Sección~\ref{sec:casos_uso}. Su función es de trazabilidad:
	permite verificar que cada caso de uso especificado tiene un módulo que lo realiza, y que ningún
	módulo se construyó sin un caso de uso que lo justifique.
	
	\rowcolors{2}{rowblue}{white}
	
	\small
	\renewcommand{\arraystretch}{1.25}
	\setlength{\tabcolsep}{4pt}
	
	\begin{longtable}{|p{4.6cm}|p{4.4cm}|p{5.0cm}|}
		\caption{Correspondencia entre módulos entregados y casos de uso}
		\label{tab:trazabilidad} \\
		\hline
		\rowcolor{headerblue}
		\headercell{Módulo} & \headercell{Casos de uso} & \headercell{Requerimientos} \\
		\hline
		\endfirsthead
	
		\hline
		\rowcolor{headerblue}
		\headercell{Módulo} & \headercell{Casos de uso} & \headercell{Requerimientos} \\
		\hline
		\endhead
	
		Registro público de quejas & CU-Q01 & RF-Q02, RF-Q03, RF-Q15 \\ \hline
		Consulta por folio & CU-Q02 & RF-Q01 \\ \hline
		Activación de cuenta y acceso & CU-Q03, CU-Q04, CU-Q05 & RF-Q04, RF-Q05, RF-Q06 \\ \hline
		Panel del quejoso & CU-Q06, CU-Q07, CU-Q09, CU-Q10 & RF-Q07 a RF-Q11, RF-Q13, RF-Q14
		\\ \hline
		Acuerdos de conciliación (quejoso) & CU-Q08 & RF-Q12 \\ \hline
		Preguntas frecuentes & --- & RF-Q16 \\ \hline
		Bandeja de entrada & CU-R02 & RF-R02 \\ \hline
		Detalle, rechazo y turnado & CU-R03, CU-R04 & RF-R03 a RF-R07, RF-R12 \\ \hline
		Registro manual & CU-R05 & RF-R08 \\ \hline
		Historial y exportación & CU-R06 & RF-R09, RF-R10 \\ \hline
		Emisión de acuerdos & CU-R07 & RF-R11 \\ \hline
		Usuarios y roles & CU-A02, CU-A03, CU-A04 & RF-A03 \\ \hline
		Catálogo institucional & CU-A05 & RF-A04 \\ \hline
		Plantillas de documentos & CU-A06 & RF-A05 \\ \hline
		Seguridad y respaldos & CU-A07, CU-A08, CU-A09 & RF-A06, RF-A07 \\ \hline
		Administración de preguntas & CU-A10 & RF-A08 \\ \hline
	\end{longtable}
	
	\normalsize
	\setlength{\tabcolsep}{6pt}
	\renewcommand{\arraystretch}{1}
